\BeleskeID{01-s013}
% element: 01-s013:e01
% element: 01-s013:e02
% element: 01-s013:e03
% element: 01-s013:e04

% element: 01-s013:d01
U digresiji oznaka \(y\) uvodi ime izlazne promenljive, pa se veza sa ulazom piše \(y=\bar{x}\). Ako je izlaz neposredno označen sa \(\bar{x}\), isti odnos već je sadržan u nazivu signala. Šeme zbog različitih oznaka ne ostvaruju različite funkcije.

Mali kružić na simbolu označava negaciju. Crta iznad promenljive, znak \(\neg\) i tilda u jednačinama predstavljaju isto logičko komplementiranje u okviru ovog predavanja.
